% Appendix to the use case "quantitative failure of trace characterisations" (usecases/trace_characterizations.tex).
% \input at the end of main.tex. Numbers: generated/trace_characterizations.tex (macros \TC...).
\section{Trace characterisations: full table and proofs}\label{app:tc}

\begin{table}[h]
\centering\scriptsize
\setlength{\tabcolsep}{1.8pt}
\begin{tabular}{rlccccccrc}
\toprule
 & inequality $L\le R$ & ref. & $V(0.1)$ & $V(0.5)$ & $V(0.9)$ & $V^\star(0.5)$ & order & slope & $n$-dep.\\
\midrule
\IfFileExists{generated/trace_characterizations.tex}{\TCRows}{\multicolumn{10}{c}{\emph{run scripts/trace\_characterizations.py}}\\}
\bottomrule
\end{tabular}
\caption[]{The inequalities and their certified search values on $M_2$ at $\delta=0.1,0.5,0.9$. $N=R$ except
$N=\varphi(P)+\varphi(Q)$ in 9 and~21, $N=(\det A)^{1/n}$ in~11 and $N=\varphi(A+B)$ in~13. $V^\star$ is the sharp
value (Propositions~\ref{prop:tcK}, \ref{prop:tcB}, \ref{prop:tcP}, \ref{prop:tcL}); $^a$ marks a value proved sharp
for $n=2$ only, $^b$ a lower bound (Propositions~\ref{prop:tcC}, \ref{prop:tcW}), $^c$ a sharp value that is
approached but not attained. ``order'': leading order of $V^\star$ as $\delta\to0$; ``slope'':
$\log_2(V(0.2)/V(0.1))$ of the search; ``$n$-dep.'': the largest $|V_n/V_2-1|$ for $n=3,4$ and $\delta\ge0.1$; for
items 7 and 13 it reflects search shortfalls; items 19 and 20 differ for odd $n$ (Proposition~\ref{prop:tcP}). In~9, $p=1$ is
the strongest member of the family $\varphi(PQ+QP)\le2\varphi((QPQ)^p)$, $p\in(0,1]$. In~7, $r=2$ and $q=1$. In~8 we
use an extension of Young's inequality ($p=q=2$) to arbitrary $A,B$, modelled on Ando's trace form. In~11, the AM--GM
analogue. In 12 and~13, $f(t)=t^{1/2}$ (13: $f=g=t^{1/2}$); in 14, $\alpha=3$; in 15 and~16, the symmetric forms; in
24, $p=2$, $q=1$; in 25, $r=1$. ``n.r.'': no violation resolved by the search
\IfFileExists{generated/trace_characterizations.tex}{(\TCMode\ run on \TCDevice)}{}.}
\label{tab:tc}
\end{table}

\paragraph{Sources and forms.} We saw several of the cited papers only as abstracts or listings; the catalogue
\texttt{usecases/trace\_characterizations\_catalogue.md} records what we saw. The following forms are checked against
the full texts. Item~4 is~\cite[Theorem~1]{petz1988} with $k=2$; there $|\varphi(A^k)|\le\varphi(|A|^k)$ for all $A$ and
one fixed $k\ge1$ characterises the trace on $M_n$ ($k=1$ is Gardner's inequality, item~1). Item~7
is~\cite[Lemma~2]{bikchentaev2011pb} with $r=2$, $p=1$. Item~10 is the case $r=0$ of~\cite[Theorems~4, 5]{hoa2010},
which also state it on $M_n$~\cite{bikchentaev2007mono,sano2006}; item~25 is the case $r=1$, on $M_n$ due
to~\cite[Theorem~6]{hoa2007}. Item~11 is~\cite[Theorem~3(iii)]{bikchentaev2020det}, $\det(A)^{1/n}\le\frac1n\varphi(A)$
for $A\ge0$ and $\varphi(1)=n$, and item~24 is~\cite[Theorem~3(vii)]{bikchentaev2020det} with $p=2$, $q=1$; it needs
$\varphi$ faithful, which holds for $W>0$. Item~12 is~\cite[Corollary~1]{tikhonov2005} with the concave
$f(t)=t^{1/2}$. Item~13 is~\cite[Lemma~3.1 and Remark~4]{hoa2013} with $f(t)=t^{1/2}$ ($s=\frac12$ in their (10)).
Items 14, 15, 20, 21 and~22 are~\cite[Theorem~2(i), (iv), (v), (vi), (vii)]{bikchentaev2011mn}, item~14 with
$\alpha=3$. Item~16 is~\cite[Lemma~1]{bikchentaev2011pb}, stated there for positive $H,K$; both sides scale by $e^c$
under $H\mapsto H+c$, $K\mapsto K+c$, so the envelope over self-adjoint arguments is the same. Items 17, 18 and~19
are~\cite[Theorem~3.4(iv), (v), (iii)]{bikchentaev2010smj} (announced in~\cite{bikchentaev2009rm}), and item~23
is~\cite[Theorem~4]{petz1988}: a positive functional on a von Neumann algebra is tracial iff it is subadditive on the
projection lattice. In the proof of~\cite[Theorem~3.4]{bikchentaev2010smj}, (iv)$\Rightarrow$(vi), the density
$\mathrm{diag}(\frac12+s,\frac12-s)$ and $p=r_{(1,1/2+t)}$, $q=r_{(1,1/2-t)}$ give
$\varphi(pqp)-\varphi(q)=2t(2s-t-4st^2)$; the paper writes $h(t)=s-t-2st^2$, which has the same sign for $t<s/2$, so
the conclusion stands (checked in the tests). The published form of item~8 is
$\varphi(|AB|)\le\frac1p\varphi(A^p)+\frac1q\varphi(B^q)$ for $A,B\ge0$~\cite{bikchentaev2005young}; the table uses the
extension to arbitrary $A,B$, and the value $V=h-1$ of Proposition~\ref{prop:tcB}(a) refers to this extension. For
the published form with $p=q=2$, an independent search (L-BFGS from 25 random starts, $n=2,3$) gives only
$V\approx\kappa-1$: $0.1546$ at $\delta=0.5$ and $1.294$ at $\delta=0.9$, against $\kappa-1=0.1547$ and $1.294$. This
is an observation, not a proof; the linear failure of item~8 is a feature of the extension.

\subsection*{Proof of Proposition~\ref{prop:tcP}}

\emph{Step 1: reduction to rank-one projections on $M_2$.} By Lemma~\ref{lem:twoproj} (see also~\cite{halmos1969}),
$\C^n=\bigoplus_jH_j$ with $\dim H_j\in\{1,2\}$, every $H_j$ invariant under $P$ and $Q$. Every operator in $L$, $R$
and $N$ (namely $PQP$, $Q$, $|P-Q|$, $P+Q$, $1+P\wedge Q$, $e^{QPQ}$, $e^P$, $\mathrm{Re}\,PQ$, $|PQ|$, $P\vee Q$)
is then block diagonal, with block $F(P_j,Q_j)$; for the lattice operations because $\mathrm{ran}\,P\cap\mathrm{ran}\,Q
=\bigoplus_j(\mathrm{ran}\,P_j\cap\mathrm{ran}\,Q_j)$, the block projections commuting with $P$ and $Q$. Hence
$L=\sum_jL_j$ etc., where $L_j$ is the same expression for $P_j,Q_j$ and the functional $\Tr(W_j\,\cdot\,)$ on
$H_j$, with the compression $W_j$ of $W$; its spectrum lies in $[b,a]$, so $\Delta(W_j)\le\delta$. On a
one-dimensional $H_j$, $P_j,Q_j\in\{0,1\}$ commute and $L_j\le R_j$ in all seven items (e.g.\ $\max(p,q)\le p+q$,
$p+q\le1+pq$, $e^{pq}\le e^p$). On a two-dimensional $H_j$, $P_j\ne Q_j$ are rank-one projections. Since $N=\sum_jN_j$
with $N_j\ge0$, and $L_j\le R_j$ wherever $N_j=0$, we get $(L-R)/N\le\max_j(L_j-R_j)/N_j$, and the ratio does not
change when $W_j$ is scaled. It therefore suffices to bound $(L-R)/N$ for $W=\frac12(1+\tau\sigma_3)$, $0\le\tau\le\delta$,
and rank-one $P=\frac12(1+\mu\cdot\sigma)$, $Q=\frac12(1+\nu\cdot\sigma)$ with unit vectors $\mu\ne\nu$ in $\mathbb R^3$
(Bloch vectors). Moreover, for fixed $\mu,\nu$ the ratio is a linear-fractional function of $\tau$ (all three
quantities are affine in $W$), with positive denominator, hence monotone on $[-\delta,\delta]$; and conjugation
with $\sigma_1$ maps $\tau$ to $-\tau$ and $(\mu,\nu)$ to reflected vectors. So the bound for $\tau=\delta$ covers all
$\tau\le\delta$, and the $M_2$ envelope bounds the envelope in every dimension. In 17, 18 and 21--23, $N$ vanishes
on blocks where $P=Q=0$, so the values below, attained (or approached) on $M_2$, are attained on $M_n$ with $P,Q$
supported on the span of $e_1,e_2$. In 19 and 20, $N$ contains $\varphi(1)$. For even $n$ take $W$ with the
eigenvalues $a,b$, each $n/2$ times, and $n/2$ copies of the $M_2$ extremiser, one in each block $\mathrm{span}
\{e_a,e_b\}$; this attains the $M_2$ value. For odd $n$ at least one one-dimensional block is left over. In 19, a
two-dimensional block contributes $L_j-R_j<\lambda_1(W_j)-\lambda_2(W_j)$ (approached) and $N_j=\Tr W_j$, and a
one-dimensional one $L_j-R_j\le0$; by Ky Fan's maximum principle the sums of the $\lambda_1(W_j)$ and of the
$\lambda_2(W_j)$ over $k'\le k=(n-1)/2$ blocks are at most the sum of the $k'$ largest and at least the sum of the $k'$
smallest eigenvalues of $W$, so $(L-R)/N<k(a-b)/(k(a+b)+b)=(n-1)\delta/(n-\delta)$, approached for the spectrum
$a^{(k)},b^{(k+1)}$ with $P=1$, $Q=0$ on the remaining eigenvector. In~20 the same configuration, with $P=Q=0$ on the
remaining eigenvector, adds $b/k$ per block to $N$, i.e.\ $(1-\delta)/(n-1)$ to $\beta$ below (normalisation
$\Tr W_j=1$); this is a lower bound, and the $M_2$ value an upper bound.
\IfFileExists{generated/trace_characterizations.tex}{For $n=3$, $\delta=0.5$ the two are $\TCExpOddHalf$ and
$\TCRefExpqpqHalf$; the search on $M_3$ finds the former.}{}

\emph{Step 2: formulas on $M_2$.} Write $\varphi(x_0+x\cdot\sigma)=x_0+\delta x_3$, $\mu\cdot\nu=\cos2\theta$
($0<\theta\le\pi/2$) and $c=\cos\theta=|\langle u,v\rangle|$ for unit vectors $u,v$ with $P=uu^*$, $Q=vv^*$. Then
$PQP=c^2P$, $(P-Q)^2=\sin^2\theta\cdot1$, $P\wedge Q=0$, $P\vee Q=1$, $\mathrm{Re}\,PQ=\frac14(1+\mu\cdot\nu+(\mu+\nu)\cdot
\sigma)$, $|PQ|=(QPQ)^{1/2}=cQ$, $e^{QPQ}=1+(e^{c^2}-1)Q$ and $e^P=1+(e-1)P$, and $|\mu+\nu|=2\cos\theta$. We use:
\emph{for $\beta>|z|$, $\max_{|\xi|=1}(\alpha+y\cdot\xi)/(\beta+z\cdot\xi)$ is the unique zero $\lambda$ of
$\alpha-\lambda\beta+|y-\lambda z|$}, since $(\alpha+y\cdot\xi)-\lambda(\beta+z\cdot\xi)\le\alpha-\lambda\beta+|y-\lambda z|$
with equality for $\xi\parallel y-\lambda z$, and the right side decreases strictly in $\lambda$.

(17) $L/R=(1+\mu\cdot\nu)(1+\delta\mu_3)/(2(1+\delta\nu_3))$. For fixed $\mu$, the lemma with $\alpha=\beta=1$, $y=\mu$,
$z=\delta e_3$ gives $\max_\nu(1+\mu\cdot\nu)/(1+\delta\nu_3)=2(1-\delta\mu_3)/(1-\delta^2)$, so
$L/R\le(1-\delta^2\mu_3^2)/(1-\delta^2)\le\kappa^2$, with equality for $\mu_3=0$ and $\nu=(\mu-\lambda\delta e_3)/(\lambda-1)$,
$\lambda=2/(1-\delta^2)$.

(18) $L/R=\sin\theta/(1+\frac\delta2(\mu_3+\nu_3))\le\sin\theta/(1-\delta\cos\theta)\le\kappa$, the maximum at
$\cos\theta=\delta$, $\mu+\nu\parallel-e_3$. (Also $\varphi(|P-Q|)\le\kappa\varphi(P+Q)$ is the triangle inequality of
Proposition~\ref{prop:tcK} with $X=P$, $Y=-Q$.)

(19) $(L-R)/N=\varphi(P+Q)-1=\frac\delta2(\mu_3+\nu_3)<\delta$, approached as $\mu,\nu\to e_3$.

(20) For fixed $g=2\theta$, the pair $(\mu_3,\nu_3)$ ranges over the image of the unit sphere under
$\xi\mapsto(\mu\cdot\xi,\nu\cdot\xi)$, a filled ellipse, and the linear-fractional ratio is maximal on its boundary,
$\mu_3=\cos x$, $\nu_3=\cos(x+g)$. With $K=e-1$ and $E=e^{\cos^2(g/2)}-1$,
\[
\frac{L-R}N=\frac{\alpha+y\cdot\xi_x}{\beta+z\cdot\xi_x},\quad \alpha=\tfrac{E-K}2,\ y=\tfrac\delta2(E\cos g-K,\,-E\sin g),
\ \beta=1+\tfrac K2,\ z=(\tfrac{\delta K}2,0),
\]
$\xi_x=(\cos x,\sin x)$; by the lemma its maximum $\lambda(\delta,g)$ is the larger root of
$(\beta^2-|z|^2)\lambda^2-2(\alpha\beta-y\cdot z)\lambda+\alpha^2-|y|^2=0$. For small $\delta$ the maximiser has
$g=2K\delta/e+O(\delta^2)$, $x=\pi/2$, and $L-R=-\frac e8g^2+\frac{K\delta}2g+O(\delta^3)$, which gives
$V=K^2\delta^2/(e(e+1))+O(\delta^3)$.

(21) With $s=\mu_3+\nu_3\ge-2\cos\theta$, $(L-R)/N=-(2\cos^2\theta+\delta s)/(4(1+\frac\delta2s))$, which decreases in $s$,
so $(L-R)/N\le(\delta C-C^2)/(2(1-\delta C))$, $C=\cos\theta$. The maximum over $C$ is at
$\delta C^2-2C+\delta=0$, $C=\tan(\psi/2)$, with value $C^2/2=\frac12\tan^2(\psi/2)$.

(22) $\varphi(\mathrm{Re}\,PQ)=\mathrm{Re}\,\varphi(PQ)\le|\varphi(PQ)|\le\kappa\varphi(|PQ|)$ by
Proposition~\ref{prop:tcK}, in every dimension. Equality: $v=\sqrt{b/(a+b)}\,e_1+\sqrt{a/(a+b)}\,e_2$, $u=Wv/\|Wv\|$;
then $PQ=\langle u,v\rangle uv^*$ has the polar factor of Proposition~\ref{prop:tcK}.

(23) $L/R=1/(1+\frac\delta2(\mu_3+\nu_3))<1/(1-\delta)$, approached as $\mu,\nu\to-e_3$. \qed

\subsection*{Proof of Proposition~\ref{prop:tcL}}
For $A\ge0$, $\varphi(A^2)=\Tr(A^{1/2}WA^{1/2}A)\le\|A\|\varphi(A)\le\Tr(A)\varphi(A)\le\varphi(A)^2/w_{\min}$, with
equality for the projection onto an eigenvector for $w_{\min}$. With $\Tr W=n$ and the extreme eigenvalues in the
ratio $h$, $w_{\min}=nb/\sum_iw_i\ge nb/(b+(n-1)a)=n\ell/(\ell+n-1)$, with equality for the spectrum
$\{a,\dots,a,b\}$. For $n=2$, $(\ell+1)/(2\ell)=1/(1-\delta)$. \qed

\subsection*{Proof of Proposition~\ref{prop:tcW}}
If $A=\alpha uu^*$ is of rank one and $A\le B$, then $u^*Bu\ge\alpha$ and $AB^2A=\alpha^2\|Bu\|^2uu^*\ge\alpha^2(u^*Bu)^2
uu^*\ge A^4$, so every positive $\varphi$ satisfies the inequality. For the construction, take $W$ with $\Tr W=1$,
eigenvalues $(1\pm\delta)/2$, and, in an orthonormal basis at the angle $\vartheta$ to its eigenbasis,
$A=\mathrm{diag}(1,a)$ and $B=A+tvv^*$ with $v=(\cos\phi,\sin\phi)$; there $W$ has the entries
$w_{11,22}=\frac12(1\pm\delta\cos2\vartheta)$ and $w_{12}=\frac\delta2\sin2\vartheta$. Then $R-L=tT_1+t^2T_2$ with
$T_1=\varphi(A^2PA+APA^2)$, $T_2=\varphi(APA)$, $P=vv^*$, and the best $t=-T_1/(2T_2)$ gives $(L-R)/R=G/(L-G)$,
$G=T_1^2/(4T_2)$, whenever $T_1<0$. Let $m_{25}(\delta)$ be the maximum over $(\vartheta,a,\phi)$; it is attained,
and by Step~1 above (applied to $\C^2\subset\C^n$) it is a lower bound in every dimension. For the leading order put
$\vartheta=\pi/4$, $\sqrt a=\sigma\delta$ and $\tan\phi=-u\,a^{-3/2}$. Then
$T_1=\cos^2\phi\,(1+u^2-(1+a)u/\sigma)$, $T_2=\cos^2\phi\,(\frac12+\frac{u^2}{2a}-\frac u\sigma)$ and
$\cos^2\phi=a^3/(a^3+u^2)$, so that
\[
\frac{L-R}R=\bigl(\sigma^2(s-s^2)-\sigma^4\bigr)^2\delta^8+O(\delta^{10}),\qquad s=\sigma/u,
\]
which is maximal for $s=\frac12$, $\sigma^2=\frac18$: $m_{25}(\delta)\ge\delta^8/4096+O(\delta^{10})$. The numerical
maximisation gives $m_{25}(0.1)/(0.1^8/4096)=\IfFileExists{generated/trace_characterizations.tex}{\TCWmonLead}{1.02}$.
The same scaling for general $r\ge0$ gives the order $\delta^{4(r+1)}$; for $r=0$ the supremum of its leading
coefficient is $\frac1{16}$, that of item~10. \qed

\paragraph{Items 12--16.} These have no closed form so far. For~14 a short expansion gives the leading order.
Write $W=\frac12(1+\delta D)$ with $D=\mathrm{diag}(1,-1)$, take real unit vectors $u,v$ at angles $m\pm\varepsilon/2$
in the eigenbasis of $W$, and set $P=uu^*$, $Q=vv^*$. Then
$L-R=\frac12(-\varepsilon^2-\delta(1-\frac2p)\varepsilon\sin2m)+O(\delta^3)$ for $\varepsilon=O(\delta)$.
Hence $V_{14}\ge(1-\frac2p)^2\delta^2/4+O(\delta^3)$, which is $\delta^2/36$ for $p=3$, while for $p=2$
every positive $\varphi$ satisfies the inequality, in agreement with the remark after~\cite[Theorem~2]{bikchentaev2011mn}
that the case $\alpha=2$ does not characterise traces. For item~13 the natural normalisation $N=R$
would give $V=\infty$ for every $\delta>0$: for rank-one $A=uu^*$, $B=vv^*$ with overlap $c\to0$,
$L\approx2c\,|u^*Wu_\perp|$ while $R=2c^2u^*Wu$. We therefore take $N=\varphi(A+B)$. The search gives:
\begin{itemize}
\item all five items fail quadratically as $\delta\to0$, with $V\approx0.5\,\delta^2$ (12, 13), $\delta^2/36$ (14),
$0.064\,\delta^2$ (15) and $0.09\,\delta^2$ (16);
\item item~15 decreases with $n$, because the identity part of $e^{P/2}e^Qe^{P/2}$ enters $N$; the others give the
same values for $n=2,3,4$, up to search shortfalls of item~13 at $\delta\ge0.9$;
\item in item~16 the best points have large $\|A\|$, $\|B\|$, and the supremum is approached as the scale grows.
\end{itemize}
An independent SciPy search on $M_2$ (plain NumPy, BFGS from
\IfFileExists{generated/trace_characterizations.tex}{\TCSciStarts}{many} random starts, at $\delta=0.3,0.5,0.9$)
gives values whose ratio to the library's lies in
\IfFileExists{generated/trace_characterizations.tex}{$[\TCSciRatioMin,\TCSciRatioMax]$}{(not generated)}. Ratios
above~1 come from item~13, where BFGS stops earlier than the batched search.

\paragraph{Items 17--25 in the search.} For the projection items the search on $M_n$, $n=2,3,4$, reproduces the values
of Proposition~\ref{prop:tcP}, including the odd-$n$ values of 19 and~20
\IfFileExists{generated/trace_characterizations.tex}{(for 17, 18 and 21--23 on $M_3$ and $M_4$ to $\TCNDevNew$ of the
$M_2$ values)}{}. In 19 and~23 the suprema are approached as $P$ and $Q$ merge, where $P\wedge Q$ and $P\vee Q$ jump;
the search therefore maximises $(L-R)/N-10^{-12}/\|P-Q\|_{\mathrm F}^2$, which keeps $P$ and $Q$ about $10^{-3}$ apart,
while the reported and certified values are those of $(L-R)/N$. The search resolves item~25 for
\IfFileExists{generated/trace_characterizations.tex}{$\delta\ge\TCWmonFirst$, reaches $\TCWmonRatioHalf$ of $m_{25}$ at
$\delta=0.5$ and $m_{25}$ itself (to $10^{-3}$) for $\delta\ge\TCWmonReach$; $m_{25}(0.9)=\TCRefWmonNine$, search
$\TCVWmonNine$.}{large $\delta$.} A BFGS search over $A,C\ge0$ on $M_2$ (tests) also reaches $m_{25}(0.9)$ and does not
exceed it.
